\begin{table}[htbp]\centering
\begin{threeparttable}
\caption{Removing the target species from the LSU reference as well.}
\label{tab:external}
\small
\begin{tabular}{llrr}
\toprule
Target & V18 & False known (v1) & With N6 (v2)\\
\midrule
\emph{G.\ chinense}$^{\dagger}$ & full & 28/35 & 5/35\\
\emph{G.\ chinense}$^{\dagger}$ & target removed & 28/35 & 5/35\\
\emph{G.\ rugosae}$^{\dagger}$ & full & 4/75 & 4/75\\
\emph{G.\ rugosae}$^{\dagger}$ & target removed & 75/75 & 62/75\\
\midrule
34 first-outcome folds & full & 19/673 (2.8\%) & 9/673 (1.3\%)\\
34 first-outcome folds & target removed & 78/673 (11.6\%) & 53/673 (7.9\%)\\
\bottomrule
\end{tabular}
\begin{tablenotes}\footnotesize
\item False-known calls among the main-type copies of each fold's target species.
$^{\dagger}$\,Development folds: these two species informed the N6 rule. \emph{Full}: the fold's
copies placed on the intact V18; \emph{target removed}: V18 realigned without the target's records, model
refitted, all copies re-placed. First-outcome folds pool the species whose V18 records were removed; 1 further fold(s) whose species V18 does not hold by label are reported apart.
\end{tablenotes}
\end{threeparttable}
\end{table}
